\section{Grundlagen}

\subsection{Quaternionen und QVNNs}

Die Quaternionen $\mathbb{H}$ erweitern die komplexen Zahlen um drei imaginäre Einheiten $\mathbf{i},\mathbf{j},\mathbf{k}$. Für diese gelten
\[
\begin{aligned}
\mathbf{i}^2&=\mathbf{j}^2=\mathbf{k}^2=-1,\\
\mathbf{i}\mathbf{j}&=\mathbf{k},\quad
\mathbf{j}\mathbf{k}=\mathbf{i},\quad
\mathbf{k}\mathbf{i}=\mathbf{j}.
\end{aligned}
\]
Vertauscht man die Reihenfolge, ändert sich das Vorzeichen. Deshalb ist das Hamilton-Produkt im Allgemeinen nicht kommutativ. In einem QVNN werden die linearen Operationen eines MLP durch quaternionwertige Operationen ersetzt. Eine Schicht kann vereinfacht als
\[
s_j^{(l)} = \phi\left(\sum_i w_{ji}^{(l)} \otimes s_i^{(l-1)} + b_j^{(l)}\right)
\]
geschrieben werden. Dabei sind $s$, $w$ und $b$ Quaternionen, $\otimes$ ist das Hamilton-Produkt und $\phi$ ist meist eine split-Aktivierung: Eine reellwertige Aktivierungsfunktion wird komponentenweise auf $q_0,q_1,q_2,q_3$ angewendet.

Der Vorteil dieser Darstellung ist die kompakte Kopplung von vier Komponenten. Der Nachteil liegt in der Optimierung. Ein Quaternion kann zwar als Vektor in $\mathbb{R}^4$ gespeichert werden, die vier Komponenten sind durch das Hamilton-Produkt aber nicht unabhängig. Konjugationen und Produktreihenfolgen bestimmen die Richtung des Updates. Dadurch können Gradientenkomponenten, die in einem reellen Netzwerk unproblematisch gemittelt werden, in $\mathbb{H}$ gegeneinander arbeiten. Die Baselines in diesem Projekt sind deshalb nicht reellwertige MLPs, sondern quaternionwertige Lernregeln aus der Literatur.

\subsection{CMLP und GHR als Baselines}

Der CMLP-Ansatz stammt aus der Clifford-Algebra-Perspektive und verwendet split Aktivierungen sowie eine komponentenweise Pseudoableitung. Für Quaternionen ergibt sich eine praktische Lernregel, bei der Fehlerterme mit konjugierten Aktivierungen bzw. Gewichten kombiniert werden \cite{buchholz2008cmlp,bill2023comparison}. Weil CMLP in der Vergleichsstudie von Bill et al. auf mehreren Regressionsaufgaben am besten abschneidet, dient CMLP hier als Hauptreferenz \cite{bill2023comparison}.

GHR steht für Generalized Hamilton-Real Calculus. Dieser Kalkül formuliert Ableitungen in $\mathbb{H}$ formaler und verwendet Ableitungen bezüglich $q$ und $q^*$ \cite{xu2015enabling}. GHR ist mathematisch reichhaltiger, aber in der Implementierung aufwendiger.

\subsection{Predictive Coding Networks}

Ein PCN besitzt latente Variablen $x^{(1)},\dots,x^{(L)}$ und eine beobachtete Eingabe $x^{(0)}$. Jede Schicht $l+1$ sagt die darunterliegende Schicht $l$ voraus:
\[
\begin{aligned}
a^{(l)} &= W^{(l)}x^{(l+1)},\\
\hat{x}^{(l)} &= f(a^{(l)}),\\
\varepsilon^{(l)} &= x^{(l)}-\hat{x}^{(l)}.
\end{aligned}
\]
Die Energie des Netzwerks ist die Summe der lokalen Vorhersagefehler:
\[
\mathcal{E}=\frac{1}{2}\sum_{l=0}^{L-1}\left\|\varepsilon^{(l)}\right\|^2.
\]
Beim Training werden zuerst die latenten Zustände durch mehrere Inferenzschritte angepasst. Danach werden die Gewichte aktualisiert. Diese Trennung ist zentral: Ein PCN kann intern eine niedrige Free Energy erreichen, ohne dass die freie Test-Inferenz bereits eine gute Vorhersage liefert. Genau dieser Unterschied zwischen interner Fehlerreduktion und Generalisierung ist in den Experimenten sichtbar.

Für supervised Learning wird ein Output-Fehler ergänzt. Während des Trainings kann der Zielwert in die Inferenz einbezogen werden. Bei der Auswertung muss die Vorhersage dagegen ohne bekannten Zielwert durch freie Inferenz entstehen \cite{stenlund2025pcn,whittington2017approximation,millidge2022predictive}. Optimierungstechnisch unterscheidet sich ein PCN damit deutlich von einem klassischen NN: Neben der Gewichtsanpassung muss auch das Inferenzverfahren stabil sein. Ein scheinbar sinkender Trainingsfehler ist deshalb nur dann aussagekräftig, wenn die gleiche Konfiguration auch ohne geklemmtes Ziel sinnvolle Testvorhersagen erzeugt.
